% ============================================================
%  MÓDULO 22 — Reprodução WSL para leigos absolutos (R100/R101)
%  Do Explorer do Windows ao PDF, passo a passo copiável
% ============================================================
\chapter{Reprodução no Windows com WSL: Do Clique ao PDF}

% R901-PONTE-ABERTURA
O mapa teórico esclareceu o alcance das referências. Vamos agora realizar uma operação delimitada: recompilar o documento existente e conservar os registros necessários à conferência.


\roteirodidatico
{Você usa Windows, nunca abriu o terminal do Linux e recebeu um caminho estranho com wsl.localhost. Este capítulo pega sua mão: do clique no Explorer à recompilação do livro, cada comando copiável, cada saída esperada e cada erro comum com conserto.}
{\emph{WSL} é o Linux dentro do Windows; \emph{Terminal} é o balcão de comandos; \emph{Clone} é trazer a casa para sua máquina; \emph{Doctor} é o check-up; \emph{PDF} é o livro recompilado por você.}
{Abra o Explorer, depois o Ubuntu, cole cada bloco na ordem e compare sua saída com a esperada antes de seguir.}
{Outro computador pode divergir por versão ou credencial; guardar a saída real é parte do método, não fracasso.}
{Qual passo travou você e que mensagem apareceu na tela?}

\casoguiado{o caminho estranho de Ana}
{Ana recebe \pth{\\\\wsl.localhost\\Ubuntu\\home\\marceloclaro\\opencode-ecosystem-core\\livro-core} e acha que é erro. Bruno explica: é o endereço da pasta do livro vista do Windows. Ana clica, vê os arquivos \pth{main.tex} e \pth{main.pdf}; Bruno abre o mesmo lugar pelo terminal e recompila. No fim, os dois olham o mesmo PDF: ela pelo clique, ele pelo comando. Os dois caminhos devem apontar para o mesmo arquivo; a reprodução ainda exige registrar fontes, ambiente, compilação e inspeção do resultado.}

\section{Passo zero: abrir a pasta pelo Windows}

\begin{metodo}
No Windows, abra o Explorer e cole na barra de endereço: \pth{\\\\wsl.localhost\\Ubuntu\\home\\marceloclaro\\opencode-ecosystem-core\\livro-core} e tecle Enter. Você deve ver arquivos como \pth{main.tex}, \pth{macros.tex} e \pth{main.pdf}. Se der erro de rede, o Ubuntu não está rodando: abra o aplicativo Ubuntu uma vez e tente de novo. Este é o único passo sem comando: só olhar.
\end{metodo}

\begin{checkpoint}
Vejo \pth{main.pdf} e consigo abri-lo? Se sim, a ponte Windows-WSL funciona. Se não, pare aqui e abra o Ubuntu antes de continuar.
\end{checkpoint}

\section{Doze micro-passos copiáveis no Ubuntu}

\elemento{Abrir o terminal e chegar na casa}{\metainfo{ID}{WSL-01} \hfill \metainfo{Tempo}{5 min}}

\begin{metodo}
Abra o aplicativo Ubuntu. Cole um bloco por vez, tecle Enter e compare com a saída esperada.
\begin{verbatim}
pwd
ls ~
\end{verbatim}
Esperado: seu usuário e pastas. Depois entre na casa:
\begin{verbatim}
cd ~/opencode-ecosystem-core
pwd
ls livro-core/main.tex
\end{verbatim}
Esperado: o caminho termina em \texttt{opencode-ecosystem-core} e o arquivo \texttt{main.tex} é listado. Erro comum: pasta não existe. Conserto: confira o nome do usuário com \texttt{whoami} e ajuste o caminho.
\end{metodo}

\elemento{Check-up do médico da casa}{\metainfo{ID}{WSL-02} \hfill \metainfo{Rota}{\pth{marceloclaro/cli.py}}}

\begin{metodo}
\begin{verbatim}
.venv/bin/python -m marceloclaro.cli doctor
\end{verbatim}
Esperado: um relatório JSON com modo, escopo, checks e estados. Nesta edição, o diagnóstico estrutural possui 22 checks; o número pode mudar. Confira cada resultado e o escopo \cod{installation_only}. Amarelo para CLI externa ausente é normal: anote como dependência opcional, não como falha total. Erro comum: \texttt{No module named marceloclaro}. Conserto: você não está na raiz da casa; volte com \texttt{cd ~/opencode-ecosystem-core}.
\end{metodo}

\elemento{Guarda de sanidade do LaTeX}{\metainfo{ID}{WSL-03} \hfill \metainfo{Rota}{\pth{livro-core/check.py}}}

\begin{metodo}
\begin{verbatim}
cd livro-core
python3 check.py
\end{verbatim}
Esperado literal: \texttt{ok: chaves, elemento, unicode e flowch limpos.} Qualquer outra saída é um problema listado por arquivo e linha: corrija antes de compilar. Erro comum: tentar compilar com erro de chaves. Conserto: não compile; leia a linha indicada.
\end{metodo}

\elemento{Recompilar o livro}{\metainfo{ID}{WSL-04} \hfill \metainfo{Tempo}{5 a 15 min}}

\begin{metodo}
\begin{verbatim}
latexmk -pdf -interaction=nonstopmode -halt-on-error main.tex
ls -lh main.pdf
\end{verbatim}
Esperado: compilação sem linha iniciada por \texttt{!} e um \texttt{main.pdf} com data de agora e tamanho de alguns MB. Abra-o pelo Explorer no caminho do passo zero e confira o sumário: devem aparecer os capítulos novos (Projeto Pedagógico, Mapa das Teorias, Reprodução WSL, Atlas). Erro comum: \texttt{pdflatex: command not found}. Conserto: instale com \texttt{sudo apt install latexmk texlive-latex-extra texlive-lang-portuguese texlive-fonts-recommended poppler-utils} e repita. Erro comum 2: travar em fonte. Conserto: aguarde; primeira compilação é lenta.
\end{metodo}

\begin{dicaleigo}
Em 30 segundos: você refez o livro. O PDF com data de hoje é sua fotografia: prova que o caminho é percorrível nesta máquina, não que toda resposta do sistema é correta.
\end{dicaleigo}

\begin{paratecnico}
Onde conferir: hash e data com \texttt{sha256sum main.pdf} e \texttt{git rev-parse HEAD} na raiz. Anote os dois ao lado da contagem de páginas. Fotografia datada é o padrão FAIR deste livro.
\end{paratecnico}

Cada entrega carrega proveniência para ser encontrada e reutilizada, segundo os princípios FAIR — encontrabilidade, acessibilidade, interoperabilidade e reutilização — \citref{WILKINSON, M. D. et al. The FAIR Guiding Principles for scientific data management and stewardship. Scientific Data, v.~3, 160018, 2016}{10.1038/sdata.2016.18}{Exige proveniência por observação em cada passo: comando, versão, hash.}{There is an urgent need to improve the infrastructure supporting the reuse of scholarly data}{Há necessidade urgente de melhorar a infraestrutura que sustenta a reutilização de dados acadêmicos.}{WILKINSON et al., 2016, 160018}.

A exigência de acompanhar o ambiente junto ao texto vem da ciência reproduzível: sem ele, o artigo descreve intenção; com ele, outro percorre o caminho \citref{PENG, R. D. Reproducible research in computational science. Science, v.~334, n.~6060, p.~1226-1227, 2011}{10.1126/science.1213847}{Justifica tratar comandos e versões como parte do texto.}{Reproducible research is the idea that the ultimate product of academic research is the paper along with the full computational environment}{Pesquisa reproduzível é a ideia de que o produto final é o artigo junto com o ambiente computacional completo.}{PENG, 2011, p.~1226-1227}.

\begin{aviso}
Limite WSL. Os recibos de 11/10/2026 pertencem à edição histórica. A compilação e a conferência desta atualização são registradas separadamente no manifesto R901. Modelos locais, credenciais de provedores e pacotes do Ubuntu mudam com o tempo. Divergência de versão não é mentira do livro: é convite a registrar nova fotografia com data, hash e saída real.
\end{aviso}

\begin{resultado}
\textbf{Exercícios.} (1) Anote commit, tamanho do PDF e saída do doctor lado a lado. (2) Provoque um erro de propósito (troque uma letra de um comando) e copie a mensagem. (3) Explique a alguém em três frases o que você provou ao recompilar e o que você não provou.
\end{resultado}

\section[Leitura guiada]{Leitura guiada: conferir a reprodução no próprio ambiente}
\begin{descricao}
\textbf{O fenômeno.} O mesmo arquivo pode produzir saídas diferentes quando versões, configuração ou recursos do host mudam.
\end{descricao}
\begin{definicao}
\textbf{Grandezas e definições.} A fonte é o conjunto de arquivos editáveis; o ambiente contém versões e executáveis; o artefato é a saída produzida. O SHA-256 identifica bytes e não representa qualidade científica.
\end{definicao}
\begin{metodo}
\textbf{Dedução passo a passo.} Abra o terminal no diretório indicado; registre a revisão e o ambiente; execute o diagnóstico; confira o código de retorno; compile; identifique o PDF por caminho e hash; abra páginas representativas. Conserve mensagens e divergências junto da data.
\end{metodo}
\begin{resultado}
\textbf{Exemplo resolvido.} Bruno recompila e encontra um hash diferente do PDF recebido. Ele verifica fontes e versões antes de concluir que há erro. Metadados de criação podem mudar os bytes; a conferência do conteúdo e do processo é distinta da igualdade de hash.
\end{resultado}
\begin{resultado}
\textbf{Ordem de grandeza e verificação.} Confira uma execução completa com seus arquivos, log e PDF. Uma resposta de diagnóstico comprova somente seu escopo; não demonstra que todos os modelos ou serviços foram executados.
\end{resultado}
\begin{aviso}
\textbf{Limite de validade.} Os resultados registrados em 11/10 pertencem àquele checkout. Reproduzi-los em outro host exige nova observação. Compilar não valida citações, eficácia científica ou comportamento de todos os providers.
\end{aviso}
\begin{resultado}
\textbf{Exercícios.} Registre ambiente e hash; selecione uma página de figura e uma de referências; explique uma divergência observada; separe uma prova de instalação de uma prova de execução.
\end{resultado}


% BEGIN-R911-ATLAS

\begin{descricao}\textbf{Navegação visual.} Infraestrutura, dados e segurança local: seção~\ref{arch:AR16}. Cada seção contém mapa mental, diagrama estrutural e fluxograma com fontes e limites.\end{descricao}

% END-R911-ATLAS

% R901-PONTE-FECHO
A compilação oferece um primeiro resultado reproduzível, mas não encerra a avaliação do Core. O epílogo reunirá o percurso completo da nota científica e as perguntas que permanecem abertas.
